% measurement_cover_letter.tex — cover letter for the Measurement submission.
% ITEMS MARKED [CONFIRM] NEED THE AUTHORS' DECISION BEFORE SUBMISSION.
% Build from paper/ with pdflatex.
\documentclass[11pt,a4paper]{letter}
\usepackage[margin=25mm]{geometry}
\usepackage[T1]{fontenc}
\usepackage{newtxtext,newtxmath}
\usepackage[hidelinks]{hyperref}
\usepackage{xurl}
\signature{Yugam Jinesh Vora\\(corresponding author, on behalf of both authors)}
\address{Department of Engineering Science\\National University of
Singapore\\\texttt{yugam@u.nus.edu}}

\begin{document}
\begin{letter}{The Editor-in-Chief\\\emph{Measurement}}
\opening{Dear Editor,}

We submit the manuscript \emph{Surrogate models for Monte Carlo uncertainty
propagation with few runs of an expensive model: specification, selection
and a valid surrogate contribution}, by Yugam Jinesh Vora (National
University of Singapore) and Wee Hoe Ng (National Metrology Centre,
A*STAR), for consideration as a research article.

\emph{Contribution to measurement science.} When the measurement model is
an expensive simulation, the Monte Carlo method of GUM Supplement~1 is
applied to a surrogate fitted to a few tens of model runs, and the error of
the surrogate becomes a component of the measurement uncertainty. Which
surrogate to use, and how to include its error in the uncertainty budget so
that the stated coverage holds, are open questions in practice. The
manuscript answers both with a paired comparison on twelve models with up
to 20 inputs, including GUM-type input distributions and a computational
fluid dynamics simulation, at budgets of 10 to 80 runs, with every
uncertainty statement checked against reference Monte Carlo evaluations of
the true models. Its main findings are:
\begin{itemize}
\item the specification of the surrogate matters more than committees or
stacking of surrogates, the type of design or adaptive designs, and from
20 runs on, choosing among a few candidates by their leave-one-out error is
practically as good as the best choice in hindsight;
\item adding the surrogate's resampled leave-one-out residuals to every
Monte Carlo trial gives GUM~S1 coverage intervals with the nominal
coverage from 10 to 80 runs, also with dependent inputs and for vector
outputs, whereas intervals from the Gaussian process predictive variance,
the usual choice, often fell short of their nominal coverage, also with the
hyperparameters treated in a fully Bayesian way;
\item the results are condensed into a procedure, applied in a worked
example with an uncertainty budget for a gas-dispersion simulation.
\end{itemize}
The work fits the scope of \emph{Measurement} in the evaluation of
measurement uncertainty and the mathematical tools for measurement, and is
directed at metrologists and engineers who evaluate uncertainty with
simulation models.

\emph{Prior presentation.} A study at a budget of 20 runs on four models
has been submitted as an extended abstract to the XXV IMEKO World Congress
(Rimini, 2027) [CONFIRM status at submission]. The manuscript extends it
substantially and has not been published or submitted elsewhere.

\emph{Data and code.} All code, the per-seed result records, the cached
simulation runs and a script that recomputes every number in the
manuscript are openly available and archived at Zenodo (DOI on the title
page). The manuscript is anonymized for double-anonymized review; the
declaration on the use of generative AI is on the title page.

[CONFIRM: optional suggested reviewers, chosen by the authors, with their
affiliations and e-mail addresses.]

The authors declare no competing interests. Thank you for considering the
manuscript.

\closing{Yours sincerely,}
\end{letter}
\end{document}
